\documentclass[master, fontsize=12pt]{diploma}

\usepackage{task}

\student{Инютин Максим Андреевич}
\faculty{№ 8 <<Компьютерные науки и прикладная математика>>}
\department{Образовательный центр Института № 8}
\speciality{01.04.02 Прикладная математика и информатика}
\profile{Продуктовая разработка и разработка программных систем}
\group{М8О-206СВ-24}

\theme{Разработка алгоритма лидарной одометрии беспилотного автомобиля, учитывающего перекрытие поля зрения}
% Дата. Оставляем пустое место для дня
\workDate{\uline{\hspace{24pt}} мая \the\year\ года}

\supervisor{Ухов Пётр Александрович}
\supervisorCredentials{кандидат технических наук, доцент, преподаватель кафедры 806}

\firstConsultant{Евсеев Дмитрий Сергеевич}
\firstConsultantCredentials{доктор физико-математических наук, профессор, профессор кафедры 806}

% \secondConsultant{Алешина Елена Константиновна}
% \secondConsultantCredentials{доктор физико-математических наук, профессор, профессор кафедры 806}

% \reviewer{Свиридов Михаил Артемьевич}
% \reviewerCredentials{доктор физико-математических наук, профессор}

\headOfDepartmentTitle{И. о. директора Образовательного центра}
\headOfDepartment{Крылов Сергей Сергеевич}



\headOfDepartmentTitle{Директор Дирекции Института № 8}

\addbibresource{main.bib}

\begin{document}
\fillTaskHeader

\textbf{1.} \makeThemeLine

\textbf{2. Срок сдачи обучающимся законченной работы} \uline{24.05.\the\year \hfill}

\textbf{3. Задание и исходные данные к работе}\\
Разработать и исследовать алгоритм лидарной одометрии беспилотного автомобиля, устойчивый к частичному перекрытию поля зрения крупным динамическим объектом.
Исходные данные: последовательность лидарных облаков и эталонные оценки движения из набора данных Boreas-RT.

\textbf{Перечень иллюстративно-графических материалов:}
\vspace{\parskip}
\begin{table}
    \begin{tabularx}{\textwidth}{|c|X|c|}\hline
    \textbf{№ п/п} & \centering\textbf{Наименование} & \textbf{Количество листов} \\ \hline
    1 & Презентация к докладу & 19 \\ \hline
    \end{tabularx}
\end{table}

% Все что выше должно уместиться на одну страницу. В крайнем случая можно переносить части вниз (или наоборот)
\vfill\clearpage
\newgeometry{left=15mm, right=30mm, top=20mm, bottom=20mm} % зеркальные поля на второй странице

\textbf{4. Перечень подлежащих разработке разделов и этапы выполнения работы}
\vspace{\parskip}
\begin{table}
    \fontsize{10pt}{12pt}\selectfont
    \begin{tabularx}{\textwidth}{|c|>{\centering\arraybackslash}X|c|c|c|}\hline
    \textbf{\makecell{№\\п/п}} & \textbf{Наименование раздела или этапа} &
    \textbf{\makecell{Трудоёмкость в\\\% от полной\\трудоёмкости\\ВКРБ}} &
    \textbf{Срок выполнения} & \textbf{Примечание} \\ \hline
    1 & Анализ предметной области и существующих методов лидарной одометрии & 20\% & 09.02.\the\year--01.03.\the\year & \\ \hline
    2 & Формализация сценария перекрытия поля зрения и требований к алгоритму & 15\% & 02.03.\the\year--15.03.\the\year & \\ \hline
    3 & Разработка алгоритма обработки лидарных кадров, взвешенного ICP и модели движения & 30\% & 16.03.\the\year--12.04.\the\year & \\ \hline
    4 & Реализация экспериментального конвейера, метрик качества и средств визуализации & 20\% & 13.04.\the\year--03.05.\the\year & \\ \hline
    5 & Проведение вычислительных экспериментов и анализ результатов & 10\% & 04.05.\the\year--17.05.\the\year & \\ \hline
    6 & Оформление выпускной квалификационной работы и презентационных материалов & 5\% & 18.05.\the\year--24.05.\the\year & \\ \hline
    \end{tabularx}
\end{table}

\textbf{5. Исходные материалы и пособия}\\ % основные работы, 25-30% от приведенного в работе списка
\fullcite{slam-cadena}\\
\fullcite{point-cloud-registration-survey-huang}\\
\fullcite{kiss-icp}\\
\fullcite{icp-learnopencv}\\
\fullcite{boreas-road-trip}\\
\fullcite{pyboreas}\\
\fullcite{habr-kalman-pictures}

\textbf{6. Дата выдачи задания} \uline{09.02.\the\year \hfill}

\makeTaskSignatures
\end{document}
